\appendix
\includepdfset{pages=-,pagecommand=\thispagestyle{plain}}

\section{Anleitung zum Aufbau der Kommunikation mit der KUKA-Zelle}
\label{sec:anleitung_kuka_kommunikation}

\subsection{Ziel}

Ziel dieser Anleitung ist der Aufbau einer TCP/IP-Verbindung zwischen
dem Labor-PC und der KUKA-Zelle. Über diese Verbindung können
Steuerbefehle aus dem auf dem Labor-PC ausgeführten Simulink-Modell an
die KUKA-Steuerung übertragen werden.

Die verwendete Implementierung basiert auf der Arbeit von Friesen
\cite{friesen2021} und verwendet insbesondere das dort beschriebene
KUKA-Skript zur Kommunikation mit der Robotersteuerung.

\subsection{Hardwareseitiger Anschluss}

Zunächst werden der Labor-PC und die KUKA-Zelle über ein
Ethernet-Kabel miteinander verbunden. Hierfür muss der dafür
vorgesehene Ethernet-Anschluss der KUKA-Steuerung verwendet werden
(\texttt{X66}). Am Labor-PC kann ein beliebiger freier Ethernet-Port
verwendet werden.

\textbf{Hinweis:} Die Verwendung des korrekten Ethernet-Anschlusses an
der KUKA-Steuerung ist erforderlich, da andernfalls keine
Netzwerkverbindung zwischen den beiden Geräten aufgebaut werden kann.

\subsection{Konfiguration der Netzwerkverbindung}

Die Netzwerkkonfiguration der KUKA-Steuerung kann über das SmartPad
überprüft werden. Hierzu wird im Menü
\texttt{Einstellungen -> Inbetriebnahme -> Netzwerkkonfiguration}
die konfigurierte IP-Adresse und Subnetzmaske angezeigt. Für die
verwendete KUKA-Zelle sind diese Werte

\begin{center}
\texttt{IP-Adresse: 192.168.41.64}\
\texttt{Subnetzmaske: 255.255.0.0}
\end{center}

hinterlegt.

Die bestehende Netzwerkkonfiguration sollte grundsätzlich nicht ohne
triftigen Grund verändert werden, da Änderungen an der
Netzwerkkonfiguration weitere Auswirkungen auf die Kommunikation der
KUKA-Steuerung haben können.

In der Netzwerkkonfiguration kann außerdem überprüft werden, welche
Ports für die TCP/IP-Kommunikation freigegeben bzw. konfiguriert sind.
Für die verwendete Kommunikationsimplementierung muss insbesondere
der Port \texttt{7000} verfügbar sein.

Anschließend wird die Netzwerkkonfiguration des Labor-PCs angepasst.
Hierzu wird der Ethernet-Adapter ausgewählt, über den die Verbindung
zur KUKA-Zelle hergestellt wurde. Dem Labor-PC wird eine IP-Adresse
im selben IP-Subnetz wie der KUKA zugewiesen. Bei Verwendung der oben
genannten Konfiguration kann beispielsweise

\begin{center}
\texttt{IP-Adresse: 192.168.41.1}\
\texttt{Subnetzmaske: 255.255.0.0}
\end{center}

verwendet werden.

Die Verbindung kann anschließend über die Eingabeaufforderung (CMD)
mit folgendem Befehl überprüft werden:

\begin{verbatim}
ping 192.168.41.64
\end{verbatim}

Erhält der Labor-PC eine Antwort von der KUKA-Steuerung, sind beide
Geräte auf IP-Ebene miteinander erreichbar. Ein erfolgreicher
Ping bestätigt jedoch noch nicht, dass die für die Anwendung
benötigte TCP-Verbindung über Port \texttt{7000} funktioniert.

\subsection{Start der KUKA-Kommunikation}

Für den Aufbau der eigentlichen Kommunikation muss zunächst
\texttt{KUKAVARPROXY.exe} auf der KUKA-Steuerung gestartet werden.
Hierfür wird am SmartPad in den Admin-Modus gewechselt.
Anschließend wird über

\begin{center}
\texttt{Einstellungen -> Inbetriebnahme -> Service -> HMI minimieren}
\end{center}

die KUKA-Anwendung minimiert. Dieser Schritt ist nur mit den
erforderlichen Administratorrechten möglich.

Nach dem Minimieren der HMI wird \texttt{KUKAVARPROXY.exe} über das
Desktop-Symbol gestartet. Das Programm stellt den TCP/IP-Server für
die Kommunikation bereit und zeigt den Verbindungsstatus des Clients
an. Für die Bedienung der Windows-Oberfläche der KUKA-Steuerung wird
eine externe Maus benötigt.

Sollte sich \texttt{KUKAVARPROXY.exe} nicht auf der KUKA-Steuerung
befinden, kann das Programm über einen USB-Stick übertragen werden.
Hierfür ist die entsprechende Anleitung von Friesen zu beachten.

Nach dem Start von \texttt{KUKAVARPROXY.exe} wird zur
KUKA-Anwendung zurückgekehrt. Anschließend wird das KRL-Programm im
Ordner \texttt{Driver/} ausgewählt und gestartet. Dieses Programm
basiert ebenfalls auf der Implementierung von Friesen und muss
während der Kommunikation aktiv sein, damit die von der externen
Anwendung empfangenen Befehle von der KUKA-Steuerung verarbeitet
werden können.

\subsection{Verbindungsaufbau mit dem Simulink-Modell}

Auf dem Labor-PC wird anschließend das entsprechende Simulink-Modell
geöffnet und die Simulation über QUARC gestartet. Für die
Kommunikation mit der KUKA-Steuerung wird im Simulink-Modell der
Block \texttt{Stream Client} verwendet. Dieser übernimmt die Rolle
des TCP-Clients, während \texttt{KUKAVARPROXY} auf der KUKA-Steuerung
als TCP-Server fungiert.

Der Ausgang \texttt{State} des Blocks \texttt{Stream Client} gibt den
aktuellen Verbindungsstatus an. Wechselt dieser auf den Wert
\texttt{2}, wurde die Verbindung zwischen Labor-PC und KUKA-Steuerung
erfolgreich hergestellt. Der Verbindungsstatus kann zusätzlich im
Fenster von \texttt{KUKAVARPROXY.exe} überprüft werden.

Optional kann anschließend der Ultraschallsensor am Greifer
angebracht werden. Im Fenster \glqq Sollwert\grqq{} kann die
gewünschte Distanz des Greifers zum Boden vorgegeben werden. Durch
Betätigen des Tasters \glqq Button\grqq{} wird die
Bewegungssequenz gestartet.

\subsection{Hinweise zur Fehlerbehebung}

Falls keine Verbindung aufgebaut werden kann, sollten zunächst die
folgenden Punkte überprüft werden:

\begin{itemize}
\item Ist das Ethernet-Kabel korrekt mit dem vorgesehenen
KUKA-Anschluss \texttt{X66} verbunden?
\item Befinden sich Labor-PC und KUKA im selben IP-Subnetz?
\item Sind IP-Adresse und Subnetzmaske korrekt eingestellt?
\item Ist die KUKA-Steuerung über \texttt{ping} erreichbar?
\item Ist der TCP-Port \texttt{7000} auf der KUKA-Steuerung
verfügbar?
\item Läuft \texttt{KUKAVARPROXY.exe}?
\item Ist das KRL-Programm im Ordner \texttt{Driver/} gestartet?
\item Läuft das Simulink-Modell über QUARC und ist der
\texttt{Stream Client} aktiv?
\end{itemize}

Vor dem Start der Kommunikations- und Bewegungssequenz ist
sicherzustellen, dass sich der Roboter in einer sicheren Situation
befindet und keine unerwartete Bewegung zu erwarten ist.
